\adnote{This chapter description from the proposal will function as the "roadmap" for the structure and contents of this dissertation chapter.}

%========================
% Chapter 4 — Results / Analysis (Proposal)
%========================

% \chapter{Results and Analysis} % uncomment if your class uses \chapter

This chapter presents results from practice-based case studies that function as research instruments for \textit{Speculocultural Technopoiesis} (ST). Each case documents what was modified, why, and with what experiential effects, and analyzes outcomes through the ST cycle and the four design territories. A cross-case synthesis concludes the chapter.

\subsubsection{Objectives}
\begin{itemize}[nosep]
	\item Present concrete tool/interaction modifications and experiential outcomes for cases.
	\item Analyze results through the ST cycle and the four domains.
	\item Extract a named pattern language and preliminary design principles.
	\item Establish limits and next steps that lead into Chapter 5 (Discussion).
\end{itemize}

\subsubsection{Intended Sections/Subsections}
\begin{itemize}[nosep]
	\item Practical Case Studies
	\begin{itemize}[nosep]
		\item Case Study 1: \textit{Buffalo Soldier}
			\begin{itemize}[nosep]
				\item Overview and Aims
				\item Context and Setup (Summary)
				\item Artifacts and Data
				\item Results (Tool and Interaction Modifications)
				\item Experiential Outcomes (Selected)
				\item Analysis via ST
				\item Limitations and Takeaways
			\end{itemize}
		\item Case Study 2: \textit{Did Sun Ra Fear Cryosleep?}
			\begin{itemize}[nosep]
				\item Overview and Aims
				\item Context and Setup (Summary)
				\item Artifacts and Data
				\item Results (Tool and Interaction Modifications)
				\item Experiential Outcomes (Selected)
				\item Analysis via ST
				\item Limitations and Takeaways
			\end{itemize}
		\item Case Study 3: \textit{Don't Play With My Hair}
			\begin{itemize}[nosep]
				\item Overview and Aims
				\item Context and Setup (Summary)
				\item Artifacts and Data
				\item Results (Tool and Interaction Modifications)
				\item Experiential Outcomes (Selected)
				\item Analysis via ST
				\item Limitations and Takeaways
			\end{itemize}
	\end{itemize}
	\item Cross-Case Synthesis
	\begin{itemize}[nosep]
		\item Pattern Language and Points of Interaction
		\item Implications for Theory and Method
		\item Limitations and Future Work
	\end{itemize}
\end{itemize}


% ----------------------------------------------------
\begin{comment}
\subsubsection{Case Study 1: \textit{Buffalo Soldier} — Cyclical Praxis}

\subsubsection{Overview and Aims}
\textbf{Aim.} Demonstrate ST’s full cycle in a spatial sonic computing environment by translating movement and marching/stepping logics into computational interaction. \\
\textbf{Design problem.} Conventional spatial audio tooling privileges uniform grids, linear transport, and vision-first control. The case explores alternatives derived from Black sonic practice (polyrhythm, call–response, groove) to restructure time, grouping, and control.

\subsubsection{Context and Setup (Summary)}
\textbf{Participants.} Small cohort of Black artist–technologists (details in Ch.~3). \\
\textbf{Site/Materials.} Studio/lab setting; motion/controller input; DAW or equivalent audio engine; custom mappings and scripts; recording of audio and screen. \\
\textbf{Procedure.} Session~0 (friction mapping) → Remixing sprints (tool/algorithm changes) → Embodying test (extended practice) → Transmitting (pattern naming, docs).

\subsubsection{Artifacts and Data}
\begin{itemize}
	\item Process: design notes, commit history, parameter diffs, patch files.
	\item Performance: multi-track audio, control streams (e.g., OSC/MIDI), screen captures.
	\item Dialogic: session transcripts, debriefs, autoethnographic memos.
	\item Transmission: named patterns, parameter schemas, short how-to examples.
\end{itemize}

\subsubsection{Results: Tool and Interaction Modifications}
\textbf{R1. Polyrhythmic time base.} Introduction of layered time lattices (e.g., 3{:}2, 5{:}4 groupings) as first-class transport, replacing a single global grid. \\
\textbf{R2. Call–response trigger graph.} Event routing that prioritizes responsive motifs (call node → response node) rather than track-centric triggering. \\
\textbf{R3. Groove-first transport.} Start/stop and seek mapped to groove anchors (downbeat clusters, turnaround points) rather than absolute bar:beat. \\
\textbf{R4. Spatial grouping by step patterns.} Spatial bus assignment and panning laws reflect stepping formations rather than Cartesian placement alone.

\subsubsection{Experiential Outcomes (Selected)}
\begin{itemize}
	\item Increased \textit{entrainment and feel}: performers reported “locking into” layered cycles more easily when transport followed groove anchors.
	\item \textit{Friction reduction}: common misfits (quantize-to-4/4; rigid clip length) were bypassed by layered lattices and call–response graphs.
	\item \textit{Technique emergence}: new micro-gestures (e.g., off-axis step accents) became reliable control primitives for sonic change.
\end{itemize}

\subsubsection{Analysis via ST}
\textbf{Unsettling.} Surfaced misfits: (i) single global grid; (ii) track-centric routing; (iii) visual-dominant transport control. \\
\textbf{Remixing.} Implemented layered time bases; response-driven routing; groove anchor transport; formation-based spatial groups. \\
\textbf{Embodying.} Cultural fit validated in extended practice; performers favored groove anchors for navigation and phrasing. \\
\textbf{Transmitting.} Named patterns: \textit{Layered Time Lattice}, \textit{Call–Response Trigger Graph}, \textit{Groove Anchor Transport}, \textit{Formation Bus Map}. Pattern writeups include parameter ranges and minimal examples.

\subsubsection{Evaluation Summary}
\textbf{Cultural authenticity.} High—interaction aligned with stepping/marching logics. \\
\textbf{Expressive affordance.} Expanded—enabled cross-cycle accents and responsive motifs. \\
\textbf{Friction reduction.} Documented—bypassed rigid quantize and track-first routing. \\
\textbf{Transmissibility.} Strong—patterns reproducible with short code and parameter schemas. \\
\textbf{Situated validity.} Positive participant uptake; interest in adopting patterns in personal workflows.

\subsubsection{Limitations / Threats to Validity}
Learning curve for layered time; room acoustics influenced spatial judgments; limited hardware profiles tested (generalization to other controllers TBD).

\subsubsection{Interim Takeaways}
\textit{Principle~1: Groove as transport.} Navigation anchored to culturally salient temporal landmarks improves phrasing and feel. \\
\textit{Principle~2: Responsivity over tracks.} Routing structures that privilege call–response better match ensemble practice than track isolation.

% ----------------------------------------------------
\subsubsection{Case Study 2: \textit{Did Sun Ra Fear Cryosleep?} — Speculative Recursion}

\subsubsection{Overview and Aims}
\textbf{Aim.} Re-apply ST to temporal structure itself: explore non-linear, elastic, and cyclic time as computational logic for composition/performance inspired by Afrofuturist temporal imagination. \\
\textbf{Design problem.} Linear bar-based transport and uniform tempo quantization limit cyclical recurrence, elastic phrasing, and return-with-difference.

\subsubsection{Context and Setup (Summary)}
\textbf{Participants.} Subset of the same cohort; emphasis on improvisatory practice. \\
\textbf{Site/Materials.} Same capture pipeline; custom temporal operators (loop orbits, stretch windows, return-with-difference transforms). \\
\textbf{Procedure.} Friction mapping of linear transport → prototyping elastic operators → extended improvisation → documentation and naming.

\subsubsection{Artifacts and Data}
\begin{itemize}
	\item Process: operator specs, code diffs, operator graphs.
	\item Performance: multi-take recordings across elastic settings; annotated return points.
	\item Dialogic: reflections on “time-feel,” recurrence, and memory.
	\item Transmission: named operators with minimal usage recipes.
\end{itemize}

\subsubsection{Results: Temporal Operators and Behaviors}
\textbf{R1. Elastic time lattice.} Local tempo regions stretch/compress within global form while preserving groove anchors. \\
\textbf{R2. Orbit loops.} Cycles that return to prior states with parameterized difference (density, timbre, voicing). \\
\textbf{R3. Memory returns.} Captured micro-phrases re-enter at culturally salient landmarks, not at fixed bar offsets. \\
\textbf{R4. Phase drift affordance.} Controlled desynchronization between layers to support polyrhythmic shimmer.

\subsubsection{Experiential Outcomes (Selected)}
\begin{itemize}
	\item \textit{Expanded phrasing}: performers sustained motifs across elastic regions without losing orientation.
	\item \textit{Recognizable recurrence}: “returns” felt intentional rather than mechanical loop restarts.
	\item \textit{Expressive drift}: slight inter-layer drift read as musical tension rather than timing error.
\end{itemize}

\subsubsection{Analysis via ST}
\textbf{Unsettling.} Identified linear transport and uniform quantization as core constraints on cyclical recurrence. \\
\textbf{Remixing.} Introduced elastic lattice, orbit loops, memory returns, and phase drift operators. \\
\textbf{Embodying.} Practice confirmed cultural salience of return points and tolerated elastic variation. \\
\textbf{Transmitting.} Named operators: \textit{Elastic Time Lattice}, \textit{Orbit Loop}, \textit{Memory Return}, \textit{Phase Drift Window}, each with parameter guidance.

\subsubsection{Evaluation Summary}
\textbf{Cultural authenticity.} High—supports cyclical time and return-with-difference in practice. \\
\textbf{Expressive affordance.} Expanded—new ways to shape density and recurrence. \\
\textbf{Friction reduction.} Demonstrated—linear transport constraints mitigated by elastic regions and memory anchoring. \\
\textbf{Transmissibility.} Moderate–strong—operators portable; requires brief onboarding to avoid disorientation. \\
\textbf{Situated validity.} Positive cohort feedback; performers reported “recognizable returns” as musically meaningful.

\subsubsection{Limitations / Threats to Validity}
Onboarding cost for elastic phrasing; risks of drift becoming noise under certain tempi; limited repertoire tested.

\subsubsection{Interim Takeaways}
\textit{Principle~3: Elastic local time.} Local modulation of time against stable groove anchors preserves orientation while expanding phrasing. \\
\textit{Principle~4: Return-with-difference.} Recurrence operators encode cultural memory without resorting to identical loops.

% ----------------------------------------------------
\subsubsection{Cross-Case Synthesis}

\subsubsection{Leverage Points and Patterns}
\textbf{Shared leverage points.} Global grid supremacy; track-first routing; visual-dominant navigation. \\
\textbf{Emergent pattern language.} 
\emph{Layered Time Lattice}, 
\emph{Call–Response Trigger Graph}, 
\emph{Groove Anchor Transport}, 
\emph{Formation Bus Map}, 
\emph{Elastic Time Lattice}, 
\emph{Orbit Loop}, 
\emph{Memory Return}, 
\emph{Phase Drift Window}. \\
\textbf{Design guidance.} Start from culturally salient anchors (groove/formation/return) rather than abstract bars and tracks.

\subsubsection{Implications for Theory and Method}
\textbf{Mediation as culturally specific.} Results concretize how defaults act as rhetorical arguments about time, grouping, and control. \\
\textbf{ST cycle robustness.} Both cases traversed Unsettling→Remixing→Embodying→Transmitting; iterative passes refined operators and pattern names. \\
\textbf{Evaluation.} Embodied cultural authenticity and transmissibility were decisive; technical correctness alone was insufficient.

\subsubsection{Limitations and Future Work}
Generalization beyond cohort requires broader repertoire and devices; some operators may require UI redesigns for clarity. Future iterations should test alternative controllers, genres, and collaborative scales.

% (Optional) Artifact/Figure Index for the Proposal Draft
% \subsection{Figures and Tables (planned)}
% Fig. 4.1 Layered Time Lattice diagram; Fig. 4.2 Call–Response Trigger Graph; 
% Fig. 4.3 Elastic Time Lattice timeline; Fig. 4.4 Orbit Loop operator graph, etc.


% --- Chapter 4 Key References (keyword-based filter) ---
\begin{refsection}[references.bib] % ensure your ch.4 items have keyword=ch4
	\nocite{*}
	\printbibliography[keyword=ch4,heading=subbibliography,title={Key References}]
\end{refsection}
\end{comment}